\section{Appendix}

% Set supplementary figure numbering
\renewcommand{\thefigure}{S\arabic{figure}}
\setcounter{figure}{0}

\subsection{Software and data availability}
\label{subsection:software_data_availability}

To maintain blinding, software and data availability in this submission are omitted, and will be included in the final version.

% \begin{itemize}

%     \item The codebase is implemented in Python using PyTorch and is available at \url{https://github.com/jkbhagatio/mini} under an MIT License for academic use.
    
%     \item Example usage of the pipeline along with replication of the paper's results can be found at \url{https://github.com/jkbhagatio/mini/notebooks}.

% \end{itemize}

\subsection{On interpretable latents}

Many neural LVM approaches focus on learning structure in a low-dimensional latent space. However, this structure can be inherently complex and difficult to interpret. NLDisco bypasses this challenge by learning individual latents directly.

\begin{figure}[h]
    \centering
    \includegraphics[width=\linewidth]{../figures/interpretable_latents_vs_latent_space.pdf}
    \caption{
        \textbf{Interpretable latents in a complex latent space.} \\
        \small This toy example highlights the utility of MINI. The two rows show two different variables (top: discrete; bottom: continuous), each uniquely encoded by the same underlying neural activity made up of three neurons' firing rates. The viridis colorbar shows the variables' values as a function of this neural activity. (\textbf{a}) Each point in the scatterplots represents a moment in time. (\textbf{b}) A projection of this activity into a 2D latent space creates tangled trajectories where variable states (e.g. 'on' and 'off' in the discrete case, and 'high' and 'low' in the continuous case) are not easily distnguished. The start and end points of the trajectories are marked by white and black dots, respectively, while arrows indicate trajectory direction. (\textbf{c}) In contrast to the tangled latent space, MINI finds individual latents corresponding to each variable, demonstrating the potential for improved interpretability of neural representations.
    }
    \label{figure:interpretable_latents_vs_latent_space}
\end{figure}

\subsection{Additional pipeline details}
\label{subsection:additional_pipeline_details}

\subsubsection{Model training}
\label{subsubsection:model_training_details}

\begin{algorithm}[h]
\caption{Model training procedure}
\label{algorithm:sded_model_training}
\begin{algorithmic}[1]
\State \textbf{Data definitions:}
\State \quad Input neural data: $\mathbf{Y} \in \mathbb{R}^{B \times S \times N_{\text{in}}}$ \textsuperscript{(Batch, Sequence length, Input neural units)}
\State \quad Target neural data: $\mathbf{Z} \in \mathbb{R}^{B \times \times N_{\text{out}}}$ \textsuperscript{(Batch, Output neural units)}
\State \textbf{Model definitions:}
\State \quad Encoder: $\mathbf{W}_{\text{enc}} \in \mathbb{R}^{D_{\text{max}} \times S \times N_{\text{in}}}$, $\mathbf{b}_{\text{enc}} \in \mathbb{R}^{D_{\text{max}}}$ \textsuperscript{(Hidden layer neurons)}
\State \quad Decoder: $\mathbf{W}_{\text{dec}} \in \mathbb{R}^{N_{\text{out}} \times D_{\text{max}}}$, $\mathbf{b}_{\text{dec}} \in \mathbb{R}^{N_{\text{out}}}$
\State \quad Transformer block (optional) $\mathbf{\theta_{t}}$: $\{\mathbf{W}_{Q,K,V} \in \mathbb{R}^{D_{\max} \times D_{\max}}, \dots\}$
\State \quad Matryoshka levels: $\mathbf{\{D_l\}_{l=1}^L}$
\State \quad Weights each level's reconstruction loss (optional): $\mathbf{\{\lambda_l\}_{l=1}^L}$
\State \quad Auxiliary loss weight: $\mathbf{\gamma}$
\\
\Procedure{train\_step}{$\mathbf{Y}, \mathbf{Z}$}    
    \Statex \Comment{\textit{--- Forward Pass ---}}
    \State $\mathbf{A} \gets \text{ReLU}(\mathbf{Y}\mathbf{W}_{\text{enc}}^T + \mathbf{b}_{\text{enc}})$ \Comment{Get encoder activations}
    \\
    \If{$S > 1$}
        \State $\mathbf{A} \gets \text{SelfAttention}(\mathbf{A}; \theta_{\text{attn}})$ \Comment{Temporally integrate over latent space}
    \EndIf
    \\
    \State $\mathcal{L}_{\text{recon}} \gets 0$
    \For{$l=1$ to $L$} \Comment{For each level...}
        \State $\hat{\mathbf{A}}_l \gets S_{\text{topk}}(\mathbf{A}_{:, :D_l})$ \Comment{Sparsify latents with batch top-$k$}
        \State $\hat{\mathbf{Z}}_l \gets \text{ReLU}(\hat{\mathbf{A}}_l \mathbf{W}_{\text{dec}, :, :D_l}^T + \mathbf{b}_{\text{dec}})$ \Comment{Reconstruct target (decoder activations)}
        \State $\mathcal{L}_{\text{recon}} \gets \mathcal{L}_{\text{recon}} + \lambda_l \cdot \text{MSLE}(\mathbf{Z}, \hat{\mathbf{Z}}_l)$ \Comment{Get reconstruction loss}
    \EndFor
    
    \Statex \Comment{\textit{--- Auxiliary Loss for Dead Latent Resurrection ---}}
    \State $\mathcal{D} \gets \text{GetDeadLatents}(\mathbf{A})$
    \If{\text{any}($\mathcal{D}$)} \Comment{Only compute if dead latents exist}
        \State $\mathbf{R} \gets \mathbf{Z} - \hat{\mathbf{Z}}_L$ \Comment{Get reconstruction residual}
        \State $\hat{\mathbf{R}} \gets \text{ReLU}((\mathbf{A} \odot \mathcal{D}) \mathbf{W}_{\text{dec}}^T + \mathbf{b}_{\text{dec}})$ \Comment{Reconstruct residual from dead latents}
        \State $\mathcal{L}_{\text{aux}} \gets \text{MSE}(\mathbf{R}, \hat{\mathbf{R}})$ \Comment{Get auxiliary loss}
    \Else
        \State $\mathcal{L}_{\text{aux}} \gets 0$
    \EndIf
    \\
    \State $\mathcal{L}_{\text{total}} \gets \mathcal{L}_{\text{recon}} + \gamma \mathcal{L}_{\text{aux}}$ \Comment{Total loss: reconstruction + auxiliary}

    \Statex \Comment{\textit{--- Backward Pass \& Parameter Update ---}}
    \State $\mathbf{g} \gets \nabla_{\theta} \mathcal{L}_{\text{total}}$ \Comment{Compute gradients, masking $\nabla\mathcal{L}_{\text{aux}}$ to dead latents}
    \State $\theta \gets \text{Update}(\theta, \mathbf{g})$ \Comment{Update model parameters}
\EndProcedure
\end{algorithmic}
\end{algorithm}

Some notes on \nameref{algorithm:sded_model_training}:

\paragraph{Dead Latent Resurrection.}
A common failure mode in training SDEDs is "latent death," where dictionary elements cease to activate for any input. We address this with an auxiliary loss designed to revive dead latents. We monitor feature activation frequencies and identify a set of dead latents $\mathcal{D}$. These latents are then trained via an auxiliary MSE loss to reconstruct the residual error of the primary model ($\mathbf{x}_{\text{target}} - \hat{\mathbf{x}}_L$). The gradients from this auxiliary loss are exclusively applied to the parameters of the dead latents, thereby encouraging them to learn useful representations without disrupting the training of active latents.

\subsubsection{Model architecture}
\label{subsubsection:model_architecture_details}

The core of our pipeline is an overcomplete sparse encoder-decoder model with a single hidden layer and ReLU activations. While we experimented with more complex architectures, such as multi-layer decoders, we found they often made interpreting the resulting latents more difficult without a significant performance gain. The Matryoshka architecture is particularly useful for datasets where features are expected to exist at multiple scales, as it helps mitigate "feature absorption" where general features are partially subsumed by more specific ones. The transformer block is most beneficial when analyzing data with meaningful temporal dynamics, as it allows the model to learn features that evolve over a sequence rather than just instantaneous neural patterns. We found that combining these components with a batch top-$k$ operator provided a robust and flexible framework that consistently outperformed L1-based sparsity penalties and yielded more interpretable features.

\subsubsection{Model evaluation}
\label{subsubsection:model_evaluation_details}

Model evaluation centers on the trade-off between reconstruction fidelity and latent dictionary sparsity. High reconstruction fidelity is trivial to achieve with a large and dense code but undermines the goal of finding disentangled, interpretable latents. NLDisco therefore first reviews dictionary quality via latent sparsity metrics (\autoref{figure:model_eval} top). The mean L0 norm measures the average number of active latents per sample, set indirectly by the batch top-$k$ parameter, while the latent activity density visualizes the firing distribution across all latents. Together these metrics reveal whether the model has learned an efficient code that uses many latents intermittently while avoiding common pitfalls such as "representational collapse" into too many constantly active latents and/or widespread "feature death", in which a large portion of the dictionary never activates. These sparsity metrics are complemented by a couple of reconstruction fidelity metrics: variance explained (R\textsuperscript{2}) and cosine similarity of reconstruction-to-target neural activity across both spatial and temporal dimensions (\autoref{figure:model_eval} bottom). High scores in these metrics confirm that the sparse code is a sufficient representation that has captured salient information present in the target neural data.

For a deeper diagnosis, several additional optional metrics are available. To check whether explanatory power is well-distributed, a variance attribution analysis quantifies the overall reconstruction variance explained by each individual latent. And to check if temporal information present in the target neural data is preserved by the latent reconstruction, a spectral frequency analysis compares the reconstruction's frequency content to that of the target neural data, given a specified duration.

\begin{figure}[h]
    \centering
    \includegraphics[width=\linewidth]{../figures/model_eval.pdf}
    \caption{
        \textbf{Model evaluation metrics.} \\
        \small In this example SAE model evaluation, a three-level Matryoshka architecture with a total of 320 latents was ised, with top-$k$ selections of 12 $\cdot B$, 24 $\cdot B$, and 36 $\cdot B$ active latents per batch size $B$(1024). The recording had a total of 200 units, and the model was trained on binned 182,868 50 ms binned spike counts. The model exhibits an approximately normal distribution for the mean latent L0 norm, few dead or tonically active latents, and high neural reconstruction fidelity as measured by the R² and cosine similarity metrics across both units and examples, all indicators of healthy training.
    }
    \label{figure:model_eval}
\end{figure}

\subsubsection{Latents evaluation}

In this paper we use prepared dashboards for viewing the activity of latents in the context of the simulated spatial navigation task and the in vivo motor control reaching task. In general when working with new datasets, we recommend spending time designing similar simple, bespoke dashboards.

\subsubsection{Hyperparameter sweeps}
\label{subsubsection:hyperparameter_sweeps}

Effective hyperparameter tuning is crucial for balancing reconstruction fidelity with latent interpretability. We typically sweep over key architectural and optimizer parameters. Architectural choices center on the Matryoshka configuration: \texttt{dsed\_topk\_map} controls the sparsity at each hierarchical level, while \texttt{dsed\_loss\_x\_map} sets the relative weight of each level's reconstruction in the total loss. Tuning these helps discover a multi-scale dictionary that matches the complexity of the neural data. We also tune \texttt{dead\_latent\_window}, which determines when an inactive latent is revived via an auxiliary residual reconstruction loss. This prevents "feature death" by giving silent latents a chance to learn useful representations. We also sweep standard optimizer parameters like learning and decay rates.

While optimal settings are data-dependent, we offer some general heuristics. The total number of latents can be initially set relative to the number of training examples (e.g.$\text{n\_examples} / 10$) and adjusted downwards based on reconstruction performance. The top-$k$ value for a given level should reflect a reasonable estimate of how many features might be co-active in a given time window. Based on our experiments, we find that a cosine learning rate schedule and the MSLE loss function often outperform static learning rates and MSE loss, respectively, and can be used as robust defaults. These targeted sweeps and informed defaults streamline the search for a model that yields a sparse, accurate, and interpretable set of latents.

\subsection{Additional results}

\subsubsection{Simulated rat spike data in a navigation task}

\begin{figure}[h]
    \centering
    \includegraphics[width=0.9\linewidth]{../figures/simulated_movement.png}
    \caption{Simulated movement along the track with the Ornstein-Uhlenbeck formula. The simulated rat sampled all the positions along the track multiple times over 600 seconds. }
    \label{figure:sim_movement}
\end{figure}

\begin{figure}[h]
    \centering
    \includegraphics[width=0.75\linewidth]{../figures/place_cells_ratemap.png}
    \caption{Ratemap of 4 simulated place cells in a 1-meter linear track. Their place field cover the whole length of the track. }
    \label{figure:ratemap_placeCells}
\end{figure}

\begin{figure}[h]
    \centering
    \includegraphics[width=0.75\linewidth]{../figures/noise_ratemap.png}
    \caption{Simulated noise along the track: the neurons had firing rates between 1 and 5 Hz, and their firing was independent of position.}
    \label{figure:ratemap_noise}
\end{figure}

\clearpage

\subsection{Methods Comparison}

Here we highlight 15 features of neural LVM methods and create a table displaying how MINI and other relevant methods compare against these features.

These features are:
\begin{itemize}
    \item \textit{Requires multimodal data}: Whether the method requires multimodal data (e.g. video data, or various forms of behavioral data, in addition to neural data).
    \item \textit{Requires trial-structured data}: Whether the method requires trial-structured data.
    \item \textit{Supports multimodal data}: Whether the method can incorporate multimodal data, even if not required.
    \item \textit{Supports trial-structured data}: Whether the method can use trial-structured data, even if not required.
    \item \textit{Learns sparse, easily identifiable latents}: Whether the method by default learns sparse, easily identifiable latents.
    \item \textit{Learns hierarchical latents}: Whether the method by default learns latents that are hierarchically organized (e.g. whether the method can learn one latent that corresponds to a particular behavior, and another, sparser latent that corresponds to a sub-behavior of the first)
    \item \textit{Learns temporally precise latents}: Whether the method by default learns latents that are time-locked to neural events at the resolution of the desired event, potentially down to single-spike precision.
    \item \textit{Learns a continuously-valued latent space}: Whether the method by default learns latents that change smoothly as a function of changes in the input neural data.
    \item \textit{Can use temporal dynamics to update the latent space}: Whether the method can use temporal dynamics (e.g. neural data or latent space history) to update the latent space.
    \item \textit{Imposes a prior on the latent space}: Whether the method imposes a predfined structure on the latent space (e.g. geometric constraints like orthogonality of latents, or distributional assumptions like a Gaussian latents).
    \item \textit{Uses nonlinear dynamics}: Whether the method learns latents from  nonlinear neural dynamics.
    \item \textit{Can be used as a generative model}: Whether the method can be used to generate new neural data samples from the learned latent space.
    \item \textit{Enforces neural data reconstruction}: Whether the method needs to perform neural data reconstruction when learning latents.
    \item \textit{Has approximate linear time scaling}: Whether the method has linear time scaling with respect to the number of data points in an example and in the dataset.
    \item \textit{Requires significant hyperparameter tuning}: Whether the method requires significant hyperparameter tuning to learn interpretable latents.
\end{itemize}

Light-green text indicates that the method has an ideal implementation of the feature, while dark-red text indicates a shortcoming.

\newcommand{\goodQual}[1]{\textcolor{green}{#1}}
\newcommand{\badQual}[1]{\textcolor{cb_red}{#1}}
\newcommand*\rot[1]{\rotatebox{90}{\parbox{3cm}{\centering\small #1}}}

\begin{table}[h]
\centering
\begin{threeparttable}
\setlength{\tabcolsep}{2.5pt}
\renewcommand{\arraystretch}{1.6}
\begin{tabular}{>{\raggedright}m{5cm}|c|c|c|c|c|c|c|c|c|c|}
\toprule
\textbf{Feature} & \rot{\textbf{MINI}*} & \rot{LangevinFlow* \cite{song_2025_langevinflow}} & \rot{CEBRA* \cite{schneider_2023_cebra}} & \rot{ST-NDT \cite{le_2022_stndt}} & \rot{AutoLFADS \cite{keshtkaran_2022_autolfads}} & \rot{UMAP** \cite{mcinnes_2018_umap}} & \rot{t-SNE** \cite{vandermaaten_2008_tsne}} & \rot{sparseNMF** \cite{hoyer_2004_sparsenmf}} & \rot{ICA \cite{comon_1994_ica}} & \rot{PCA \cite{hotelling_1933_pca}} \\
\midrule
Requires multimodal data & \goodQual{No} & \goodQual{No} & \goodQual{No} & \goodQual{No} & \goodQual{No} & \goodQual{No} & \goodQual{No} & \goodQual{No} & \goodQual{No} & \goodQual{No} \\
\hline
Requires trial-structured data & \goodQual{No} & \goodQual{No} & \goodQual{No} & \goodQual{No} & \goodQual{No\textsuperscript{\textasciicircum}} & \goodQual{No} & \goodQual{No} & \goodQual{No} & \goodQual{No} & \goodQual{No} \\
\hline
Supports multimodal data & \badQual{No\textsuperscript{\dag}} & \badQual{No} & \goodQual{Yes} & \badQual{No} & \badQual{No} & \badQual{No} & \badQual{No} & \badQual{No} & \badQual{No} & \badQual{No} \\
\hline
Supports trial-structured data & \goodQual{Yes} & \goodQual{Yes} & \goodQual{Yes} & \goodQual{Yes} & \goodQual{Yes} & \goodQual{Yes} & \goodQual{Yes} & \goodQual{Yes} & \goodQual{Yes} & \goodQual{Yes} \\
\hline
Learns sparse, easily identifiable latents & \goodQual{Yes} & \badQual{No} & \badQual{No} & \badQual{No} & \badQual{No} & \badQual{No} & \badQual{No} & \goodQual{Yes} & \badQual{No} & \badQual{No} \\
\hline
Learns hierarchical latents & \goodQual{Yes} & \badQual{No} & \badQual{No} & \badQual{No} & \badQual{No} & \badQual{No} & \badQual{No} & \badQual{No} & \badQual{No} & \badQual{No} \\
\hline
Learns temporally precise latents & \goodQual{Yes} & \goodQual{Yes} & \goodQual{Yes} & \goodQual{Yes} & \goodQual{Yes} & \badQual{No} & \badQual{No} & \goodQual{Yes} & \badQual{No} & \badQual{No} \\
\hline
Learns a continuously-valued latent space & \badQual{No\textsuperscript{\ddag}} & \goodQual{Yes} & \goodQual{Yes} & \goodQual{Yes} & \goodQual{Yes} & \goodQual{Yes} & \goodQual{Yes} & \badQual{No\textsuperscript{\ddag}} & \goodQual{Yes} & \goodQual{Yes} \\
\hline
Can use temporal dynamics to update the latent space & \goodQual{Yes} & \goodQual{Yes} & \goodQual{Yes} & \goodQual{Yes} & \goodQual{Yes} & \badQual{No} & \badQual{No} & \badQual{No} & \badQual{No} & \badQual{No} \\
\hline
Imposes a prior on the latent space & \goodQual{No} & \badQual{Yes} & \goodQual{No} & \goodQual{No} & \badQual{Yes} & \goodQual{No} & \goodQual{No} & \goodQual{No} & \badQual{Yes} & \badQual{Yes} \\
\hline
Uses nonlinear dynamics & \goodQual{Yes} & \goodQual{Yes} & \goodQual{Yes} & \goodQual{Yes} & \goodQual{Yes} & \badQual{No} & \badQual{No} & \badQual{No} & \badQual{No} & \badQual{No} \\
\hline
Can be used as a generative model & \goodQual{Yes} & \goodQual{Yes} & \goodQual{Yes} & \goodQual{Yes} & \goodQual{Yes} & \badQual{No} & \badQual{No} & \badQual{No} & \badQual{No} & \badQual{No} \\
\hline
Enforces neural data reconstruction & \badQual{Yes} & \badQual{Yes} & \goodQual{No} & \badQual{Yes} & \badQual{Yes} & \goodQual{No} & \goodQual{No} & \badQual{Yes} & \goodQual{No} & \goodQual{No} \\
\hline
Has approximate linear time scaling & \goodQual{Yes\textsuperscript{\#}} & \badQual{No} & \goodQual{Yes} & \badQual{No} & \badQual{No} & \goodQual{Yes} & \badQual{No} & \badQual{No} & \goodQual{Yes} & \goodQual{Yes} \\
\hline
Requires significant hyperparameter tuning & \badQual{Yes} & \badQual{Yes} & \badQual{Yes} & \badQual{Yes} & \badQual{Yes} & \badQual{Yes} & \badQual{Yes} & \badQual{Yes} & \goodQual{No} & \goodQual{No} \\
\bottomrule
\end{tabular}
\caption{\centering Comparison of MINI with other neural LVM methods.}
\label{table:method_comparisons}
\begin{tablenotes}[flushleft]
\footnotesize
\item *: Latent comparison visualized in \nameref{section:results}
\item **: Latent comparison visualized in \nameref{subsubsection:latent_comparisons_traditional}
\item \textsuperscript{\dag}: Not a fundamental limitation -- could be added as a feature
\item \textsuperscript{\ddag}: Enforced sparsity can cause step-like jumps in the data-to-latents mapping
\item \textsuperscript{\#}: In the standard implementation, without a transformer layer in the decoder
\item \textsuperscript{\textasciicircum}: Can bin data to create ``pseudo''-trials
\end{tablenotes}
\end{threeparttable}
\end{table}